\section{Excluding the linear-budget jump}
\label{sec:jump-to-large-bias-only-at-linear-budget}

\begin{proposition}[Exclusion of the linear-budget jump]\label{prop:linear-budget-jump-exclusion}
In the stabilized contradiction sequence, it is impossible that both orientations have \(H^\dagger\ge1\) while one has \(H_L^\dagger=0\).
\end{proposition}

\begin{proof}\leavevmode\par
\paragraph{The slice experiment.}
Suppose both orientations have \(H^\dagger\ge1\), but one has \(H_L^\dagger=0\); denote that orientation by \(X,Y\). We will use a biased patch inside a small cube slice to leave large first-side common-neighbor sets. Independent labels from the other slices will then pay for the remaining cube edges. The compatibility argument below ensures that those outer labels do not consume almost all of an internal set.

Fix cluster-exclusion parameters \(\zeta=1/100\) and rational \(\delta\) as in Proposition~\ref{prop:full-dimensional-cluster-exclusion}, and put
\(s_c=\lceil e^{n^\zeta}\rceil\), \(\upsilon=\delta/4\).
In the reversed orientation, \(H^\dagger\ge1\) permits a rational \(x_0>0\) with
\(H_{\rm rev}(x_0,1-\upsilon/4)>.95\).
Indeed the near-1 second-budget limit exceeds \(.95\) for a sufficiently small first argument, and monotonicity gives the same lower bound at \(1-\upsilon/4\). Stabilized absence at exponent \(.95\), followed by reversing the sides, gives
\begin{equation}
 |d_G(\sigma,\pi)-1/2|\le b_*:=n^{-1+.05}
 \quad\text{for}\quad
 \operatorname{width}(\sigma)\le n^{1-\upsilon/4},\quad
 \operatorname{width}(\pi)\le n^{x_0},
 \label{eq:source-6}
\end{equation}
on \(X,Y\), in either color.

Set \(g=n^{-.01}\). Since \(H_L^\dagger=0\), monotonicity and nonnegativity give \(H_L(.01,.01)=0\). Choose an available exponent \(h_0<.01\). Its surplus \(n^{-h_0}\) is much larger than \(g\). After fixing a color along a subsequence, the columns of degree at least \(1/2+2g\) have second-law mass at least \(c n^{-h_0}\). Conditioning on them costs only \(h_0\log n+O(1)\) width. Lemma~\ref{lem:stabilization} therefore supplies an available finite menu of tags \(i\), all of one color \(G\), such that
\[
 \operatorname{width}(\mu_i)\le n^{.01},\qquad
 \operatorname{width}(\nu_i)\le .011n,\qquad
 d_G(\mu_i,y)\ge1/2+2g\quad(y\in\operatorname{supp}\nu_i).
\]
All supports lie in the retained sets.

\paragraph{One slice and its odd rows.}
Split the cube into inner slices \(Q_h\), where \(h=\lfloor n^{.1}\rfloor\), indexed by outer words \(s\in Q_{n-h}\). Put \(k=\lceil n^{.2}\rceil\), and first consider one slice with fixed tag \(i\). Define the hit-conditioned first law
\[
 \rho_y(w)=\frac{\mu_i(w)1[w\sim_Gy]}{d_G(\mu_i,y)}.
\]
Provisionally draw \(Y_0\sim\nu_i\). Given \(Y_0\), draw the tuples \(W_v\sim\rho_{Y_0}^{\otimes k}\) independently at its even roles. At an odd role \(b\), the observed data are the tuples at its \(h\) internal even neighbors. Their likelihood density, relative to independent \(\mu_i\)-entries, is
\[
 L_b(y)=\prod_{\text{observed entries }w}
           \frac{1[w\sim_Gy]}{d_G(\mu_i,y)}.
\]
Let \(Z_b=\int L_b\,d\nu_i\), and let \(Z_{b,-v}\) be the corresponding predictive density with one neighboring tuple omitted. Require positivity and
\[
 Z_b\ge e^{-kh},\qquad Z_b/Z_{b,-v}\ge e^{-.2gk}
          \quad\text{for each internal }v\sim b.
\]
On passing, set \(p_b^W=\nu_iL_b/Z_b\); otherwise set \(p_b^W=\nu_i\).

Lemma~\ref{lem:gated-posterior} applies first to all \(kh\) entries, using reference \(\mu_i^{\otimes kh}\). For a deletion test, fix the other entries and use their posterior on \(Y_0\) as the prior; the next tuple has reference \(\mu_i^{\otimes k}\), and predictive density \(Z_b/Z_{b,-v}\). Hence
\[
 \Pr\{b\text{ fails a test}\}
       \le e^{-kh}+h e^{-.2gk}=\exp(-\Omega(gk)).
\]
The row laws always have width at most \(.02n\). On a passing row,
\[
 \log(N\max p_b^W)
 \le .011n+kh\bigl(1-\log(1/2+2g)\bigr)
 \le .02n,
\]
and the fallback has the smaller width \(.011n\).

For each fixed \(y_0\), the conditional test-failure probability is the same at every odd role by translation symmetry of the slice. Averaging over \(Y_0\) therefore selects one deterministic \(y_0\) for this tag for which the per-row bound holds. Fix it from now on and sample independent \(W_v\sim\rho_{y_0}^{\otimes k}\). The likelihood rule \(L_b\) is unchanged. This selection gives the same bound for each row, not simultaneous success of all rows at this stage. Conditional on \(W\), the slice reference law draws the odd labels independently from the rows \(p_b^W\).

\paragraph{The first-side common-neighbor set.}
At an even \(v\), let \(\sigma_v\) be uniform on the labels of \(\operatorname{supp}\mu_i\) hitting all its actual internal odd neighbors, if this set has size at least
\(N\exp(- (\log2-.5g)h)\); otherwise let \(\sigma_v=0\).
We claim that \(\sigma_v\) has mass 1 except with probability \(\exp(-\Omega(gk))\) in the slice reference law.

Fix every tuple except \(W_v\), and vary \(w=W_v\) under prior \(\rho_{y_0}^{\otimes k}\). For a vector \(t\) of internal odd outputs, let \(F_w(t)\) be the product of their row masses, multiplied by the indicator that all their tuple tests pass. Recompute both rows and tests at \((W_{-v},w)\). Let \(Q(t)\) be the product of the deletion-posterior rows, with fixed \(w\)-independent fallbacks if a deletion denominator is zero. Each neighboring row costs at most
\((1/2+2g)^{-k}e^{.2gk}\), so on the gate
\[
 F_w(t)\le e^{(\log2-2g)kh}Q(t).
\]
Put \(m(t)=\int F_w(t)\,d\rho_{y_0}^{\otimes k}(w)\). The true-gated event
\(m=0\) or \(m<e^{-.5gkh}Q\)
has probability at most \(e^{-.5gkh}\). Outside it, the posterior on \(w\) has atom cap
\[
 N^{-k}\exp\bigl((\log2-1.5g)kh+k n^{.01}+O(k)\bigr)
       \le N^{-k}e^{(\log2-g)kh},
\]
where \(n^{.01}=o(gh)\).

If the common-neighbor set has size \(T_0\), this posterior is supported on at most \(T_0^k\) tuples. Its total mass 1 forces
\[
 T_0\ge N e^{-(\log2-g)h},
\]
which exceeds the cutoff defining \(\sigma_v\). Adding the \(h\) possible neighboring row-test failures proves the claim.

\subsection{Compatibility}
\label{sec:compatibility}

Define the tag's two mean rows by
\[
 \pi_i=\mathbb E_W p_b^W,\qquad
 \alpha_i=\mathbb E_{W,\,Y_{N_{\rm in}(v)}\mid W}\sigma_v.
\]
Here \(N_{\rm in}(v)\) is the set of internal odd neighbors of \(v\). Their outputs in the second expectation are independent with rows \(p_b^W\). Thus \(\pi_i\) is a probability law, whereas \(\alpha_i\) can be a subprobability. Translation symmetry makes these means independent of the chosen role of the indicated parity, also when the outer word reverses the slice's internal parity. They are supported on the two sides of the tag, and \(\pi_i\) has width at most \(.02n\).

The next two auxiliary lemmas use the slice experiment, patch menu, and parameter choices fixed in this proof.

\begin{lemma}[A compatible balanced profile]\label{lem:compatible-balanced-profile}
There is one profile \(p\) for independent tag choices with
\begin{equation}
 \pi=\sum p_i\pi_i\le K/N,\qquad \sum p_i\alpha_i\le K/N,
 \label{eq:source-7}
\end{equation}
where \(K\) is constant, and the following compatibilities. Write \(f_x(y)=2\mathbf1[x\sim_G y]-1\), \(m_x=\mathbb E_\pi f_x\), \(K_\pi(x,x')=\mathbb E_\pi f_x f_{x'}\); fix \(\eta=10^{-8}\). For every tag used by \(p\):
\begin{itemize}
\item on \(\operatorname{supp}\mu_i\), \(|m_x|\le4b_*\) and \(p\{i':d_G(x;\pi_{i'})>.8\}\le\eta\);
\item on that support there is no clique of size \(s_c\) of pairs of distinct labels with \(K_\pi(x,x')>8n^{-\delta}\).
\end{itemize}
\end{lemma}

\begin{proof}
Fix \(K\) sufficiently large and consider the compact convex set of profiles satisfying \eqref{eq:source-7}. It is nonempty by availability and price separation. For an input profile in this set, let \(\pi=\sum p_i\pi_i\). We will discard \(o(N)\) first labels so that every patch supported off the discards meets all compatibility requirements for this input.

First discard labels with \(|m_x|>4b_*\). Their number is \(o(N)\) by \eqref{eq:source-6}: the uniform law on a larger signed exception, tested against the broad law \(\pi\), would contradict that estimate. On the remaining labels greedily remove disjoint cliques of size \(s_c\) for the graph \(K_\pi(x,x')>8n^{-\delta}\). Let \(V\) be the union removed. If \(|V|\ge Ne^{-n^\delta}\), use \(\pi\) as a first law in the reversed orientation and the uniform laws on these disjoint cliques as clusters, weighted in proportion to their sizes. Their aggregate is uniform on \(V\). Its width is at most \(n^\delta\), the width of \(\pi\) is at most \(\log K\), and each cluster atom is at most \(1/s_c\le e^{-n^\zeta}\).

For distinct members of a clique,
\[
 d_G(\pi;x,x')
   =\frac{1+m_x+m_{x'}+K_\pi(x,x')}{4}
   \ge\frac14+n^{-\delta}
\]
for large \(n\), since \(b_*=o(n^{-\delta})\). A diagonal pair has codegree \((1+m_x)/2\) and also meets this bound. This is a witness forbidden by Corollary~\ref{cor:eventual-cluster-absence}. Therefore \(|V|<Ne^{-n^\delta}\).

\paragraph{Excluding too many high-degree tags.}
Before clique removal, let \(U\) be the good-degree labels violating the \(\eta\)-condition. Suppose \(|U|\ge Ne^{-n^\delta}\), and let \(\mu_U\) be uniform on \(U\). Write \(D_i(x)=d_G(x;\pi_i)\). For every \(x\in U\), the mean of \(D_i(x)\) over tags is \(1/2+O(b_*)\), while tag mass greater than \(\eta\) has \(D_i(x)>.8\). Hence
\[
 \mathbb E_iD_i(x)^2
   =(\mathbb E_iD_i(x))^2+\operatorname{Var}_iD_i(x)
   \ge\frac14+c\eta
\]
for an absolute \(c>0\). Integrate this inequality over \(\mu_U\).

In \(L^2(\mu_U)\), put \(v_y(x)=1[x\sim_G y]\) and \(v_i=\mathbb E_{y\sim\pi_i}v_y\). Then
\(\|v_i\|_2^2=\int D_i(x)^2\,d\mu_U(x)\).
A fixed positive fraction of the tags therefore have \(\|v_i\|_2>1/2+c'\), for a fixed \(c'>0\). For each such tag, a fixed positive \(\pi_i\)-mass of labels has projection on \(v_i/\|v_i\|_2\) at least \(1/2+c''\), with fixed \(c''>0\). Let \(\rho_i\le C\pi_i\) be \(\pi_i\) conditioned on these labels.

Join two of these labels when their codegree under \(\mu_U\) is at least \(1/4+c'''\), with fixed \(c'''>0\) sufficiently small. This graph has bounded independence number. Indeed, for an independent set of size \(t_0\), the average of its adjacency vectors has projection at least \(1/2+c''\), but squared norm at most
\[
 1/t_0+1/4+c'''.
\]
These are incompatible for large fixed \(t_0\). The diagonal codegrees also exceed the threshold because the squared projection of each vector is at least \((1/2+c'')^2\).

Draw \(r_0=\lfloor e^{.03n}\rfloor\) labels from \(\rho_i\), conditioned to be distinct. Since \(\max\rho_i\le C e^{.02n}/N\), the collision probability before conditioning is
\[
 O(r_0^2e^{.02n}/N)=o(1).
\]
The elementary Ramsey binomial bound with the fixed independence threshold permits greedy packing of all but \(s_c^{O(1)}\) sampled labels into cliques of size \(s_c\). A positive fraction of the sample, in fact \(1-o(1)\), is retained. Make each clique a uniform cluster and weight these clusters in proportion to their sizes. If \(R_i\) is the retained sample, its aggregate obeys
\[
 U(R_i)\le\frac2{r_0}\sum_{j\le r_0}\delta_{Y_j}.
\]
Average over the distinct sample and over the selected positive-mass set of tags. The conditional sampling and tag restriction each cost only a constant in this inequality, so the final aggregate is bounded by \(O(1)\pi\).

We have obtained a cluster witness with first law \(\mu_U\) of width at most \(n^\delta\), aggregate second law of width \(O(1)\), atom cap \(1/s_c\), and codegree at least \(1/4+c'''\ge1/4+n^{-\delta}\). Corollary~\ref{cor:eventual-cluster-absence} again gives a contradiction. Thus the labels violating the high-degree-tag condition also number less than \(Ne^{-n^\delta}\).

It follows that outside \(o(N)\) discards all the input-dependent requirements for output tags can be met along with arbitrary price exclusions by availability. Separation gives an output profile satisfying the same balance constraints \eqref{eq:source-7}, increasing their fixed \(K\) at the outset as needed for discard slack. Require it supported only on tags compatible with the \textbf{input} as listed. These responses are convex and have closed graph (cliques used strict inequalities), so Kakutani proves the claim. To make the domain choice explicit, fix the balance constant large enough for both separation arguments before forming the compact profile domain. The single-tag laws and support sets are fixed throughout. For a fixed label, the exceptional-tag set in the degree condition is therefore fixed, and its profile mass is linear in the input profile. For each fixed candidate clique, exclusion of the strict simultaneous inequalities is a closed condition. There are finitely many labels, tags and candidate cliques at each stage, so compatibility of any fixed output tag is closed in the input profile. Restricting output mass to these compatible tags gives precisely the closed response graph used above.
\end{proof}

\subsection{Outer moment}
\label{sec:outer-moment}

\begin{lemma}[Outer mass lower tail]\label{lem:jump-outer-mass-tail}
Fix a compatible tag and any probability \(\sigma\) on its first support with
\[
 \max\sigma\le N^{-1}\exp((\log2-.5g)h).
\]
Take \(d=n-h\) independent labels \(Y_l\sim\pi\). Write \(D_x=d_G(x;\pi)\) and
\[
 a_x(y)=\frac{\mathbf1[x\sim_G y]}{D_x}-1=\frac{f_x(y)-m_x}{1+m_x}.
\]
Then for any fixed \(P\),
\begin{equation}
 \Pr\left\{\int \prod_{l\le d}(1+a_x(Y_l))\,d\sigma(x)<1/2\right\}
 \le n^{-P}
 \label{eq:source-8}
\end{equation}
eventually, uniformly.
\end{lemma}

\begin{proof}
Let
\[
 Z(Y_1,\ldots,Y_d)=\int\prod_{l\le d}(1+a_x(Y_l))\,d\sigma(x).
\]
We will bound a high even moment of \(Z-1\). For independent \(x_1,\ldots,x_u\sim\sigma\), where \(u\) is a fixed length to be chosen, define
\[
 M_J=\int\prod_{j\in J}a_{x_j}\,d\pi.
\]
Singleton interactions vanish because \(\int a_x\,d\pi=0\). A free label below means one involved tuple coordinate integrated with its original \(\sigma\) law, the other specified coordinates being fixed. All constants may depend on \(u\).

\paragraph{One free label.}
With one label free, \(|M_J|\le O_u(b_*)\) except with probability at most
\(2e^{-n^{1-\upsilon/4}/2}\).
To see this, write the interaction as
\[
 M_J=\frac{\langle f_x,H_0-\mathbb E_\pi H_0\rangle_\pi}{1+m_x},
\]
where \(H_0\) is the product of the fixed factors and is bounded by a constant depending on \(u\). If one signed exceptional set had \(\sigma\)-mass at least \(e^{-n^{1-\upsilon/4}/2}\), conditioning on it would give a first law of width at most
\(O(h)+n^{1-\upsilon/4}/2<n^{1-\upsilon/4}\).
Reweight that law by \((1+m_x)^{-1}\) and normalize, at width cost \(O(1)\). The centered second function is a difference of two probability tests obtained by adding a fixed positive offset to its positive and negative parts. Their widths are \(O(1)<n^{x_0}\). Applying \eqref{eq:source-6} to them contradicts an interaction larger than a sufficiently large constant times \(b_*\).

\paragraph{Two free labels.}
With two involved labels \(x,z\) free, we claim
\[
 \Pr\{|M_J|>n^{-1-.03}\}\le e^{-n^{.4}}.
\]
Write \(M_J=\langle a_x,H_0a_z\rangle_\pi\), with fixed bounded \(H_0\). If the claim fails, choose one sign and a set \(A\) of \(x\)'s of mass at least \(\rho=e^{-O(n^{.4})}\) such that the signed violation set \(E_x\) of \(z\)'s also has mass at least \(\rho\) for every \(x\in A\). For a fixed \(x\), draw \(t=\lceil n^{.25}\rceil\) independent \(z_j\)'s from \(\rho_x=\sigma|_{E_x}\).

The width of \(\rho_x\) is \(O(h+n^{.4})\). After the bounded reweighting \((1+m_z)^{-1}\), \eqref{eq:source-6} says that its mean adjacency indicator is within \(O(b_*)\) of \(1/2\) outside a \(\pi\)-set of mass \(e^{-\Omega(n^{x_0})}\). Otherwise that exceptional second set, conditioned with the corresponding sign, would still fit the second budget. Boundedness and the degree estimates on the first support therefore give
\[
 \|\mathbb E_{\rho_x}a_z\|_{L^2(\pi)}=O(b_*).
\]
For \(S=H_0\sum_{j\le t}a_{z_j}\), independence yields
\[
 \mathbb E\|S\|_2^2=O_u(t+t^2b_*^2)=O_u(t).
\]
A fixed positive fraction of these conditioned tuples thus has \(\|S\|_2\le C_u\sqrt t\).

Let \(\mathcal G\) denote that norm event. It depends only on the tuple and \(H_0\), not on \(x\). Under the raw product law \(\sigma^{\otimes t}\),
\[
 \int_A\Pr(\mathcal G,\ z_j\in E_x\text{ for every }j)\,d\sigma(x)
       \ge c_u\rho^{t+1}.
\]
Fubini supplies a fixed norm-good tuple whose simultaneous signed-violation set has \(x\)-mass at least \(e^{-O(tn^{.4})}\). Let \(q\) be \(\sigma\) restricted to that set. Its width is
\(O(h+tn^{.4})=O(n^{.65})<n^{1-\upsilon/4}\).

For this fixed tuple put
\[
 r_x=(1+m_x)^{-1},\qquad c_q=\mathbb E_qr_x,\qquad
 \widetilde q=r_xq/c_q,\qquad T=S-\mathbb E_\pi S.
\]
The exact projection identity is
\[
 \mathbb E_q\langle a_x,S\rangle_\pi
       =c_q\mathbb E_{\widetilde q}\langle f_x,T\rangle_\pi.
\]
The bounded summands give \(\|T\|_\infty=O_u(t)\), while centering preserves \(\|T\|_2=O_u(\sqrt t)\). Thus the positive test laws
\[
 \pi_\pm=\frac{\sqrt t+T_\pm}{Z_T}\pi,\qquad
 Z_T=\sqrt t+\mathbb E_\pi T_+
     =\sqrt t+\mathbb E_\pi T_-=O_u(\sqrt t)
\]
have equal normalizers and widths \(O_u(\log n)\). Applying \eqref{eq:source-6} to \(\widetilde q,\pi_\pm\) bounds the projection by
\(O_u(b_*\sqrt t)=O_u(n^{-.825})\).
But the simultaneous violations give magnitude at least
\(t n^{-1-.03}=n^{-.78+o(1)}\)
in one sign. This contradiction proves the two-free estimate.

\paragraph{Mean interaction size.}\label{par:jump-mean-interaction}
\(\mathbb E |M_J|\le n^{-.4|J|}\) for large \(n\).

For \(r_J=|J|\), independence and Fubini give
\[
 \mathbb E M_J^2
 =\mathbb E_{y,y'\sim\pi}
     \left(\int a_x(y)a_x(y')\,d\sigma(x)\right)^{r_J}.
\]
Fix \(y\). The inner kernel is a bounded signed reweighting of \(\sigma\) tested against \(f_x(y')\), plus an \(O(b_*)\) term arising from \(m_x\). Adding positive offsets to this signed reweighting costs only \(O(1)\) first width. By \eqref{eq:source-6}, the kernel is \(O(b_*)\) except on a \(y'\)-set of \(\pi\)-mass \(e^{-\Omega(n^{x_0})}\), uniformly in \(y\). Hence
\[
 \mathbb E M_J^2\le(C_u b_*)^{r_J}
                       +C_u e^{-\Omega(n^{x_0})}.
\]
Cauchy--Schwarz, with \(.95/2>.4\), gives the stated first absolute moment.

\paragraph{Counting large extensions.}
Fix the other labels of an interaction. There are at most \(e^{n^{.03}}\) labels in the support extending it to \(|M_J|>n^{-\upsilon}\). Such an extension has projection of \(f_x\) on a fixed bounded direction of magnitude at least \(c_un^{-\upsilon}\). Split by sign. For \(t_0\) candidates of one sign with mutual \(K_\pi\le8n^{-\delta}\), the squared norm of their average is at most
\[
 1/t_0+8n^{-\delta},
\]
whereas its squared projection is at least \(c_u^2n^{-2\upsilon}\). As \(\delta=4\upsilon\), this forces \(t_0=O_u(n^{2\upsilon})\). The compatibility hypothesis forbids an \(s_c\)-clique of larger inner products. The elementary Ramsey bound obtained from the one-vertex recursion therefore bounds the number of candidates by a binomial coefficient, with logarithm at most
\[
 \log\binom{s_c+t_0}{t_0}
       =O_u(n^{\zeta+2\upsilon})<n^{.03}.
\]
The sign split and the fixed number of interaction choices fit the slack in this bound.

\paragraph{Exponential weights for moderate interactions.}
For a tuple let \(w=\max_{|J|\ge2}|M_J|\), taking an empty maximum to be zero. On \(w\le n^{-\upsilon}\), the integral of \(e^{2^unw}\) against the original product tuple law is bounded. Its contribution from \(w>n^{-1-.03}\) is smaller than any fixed negative power of \(n\). Indeed the two ranges are bounded respectively by
\[
 \begin{aligned}
 \int_{\{n^{-1-.03}<w\le n^{-1+.06}\}}e^{2^unw}\,d\sigma^{\otimes u}
   &\le C_u e^{2^u n^{.06}-n^{.4}},\\
 \int_{\{n^{-1+.06}<w\le n^{-\upsilon}\}}e^{2^unw}\,d\sigma^{\otimes u}
   &\le C_u e^{2^u n^{1-\upsilon}-n^{1-\upsilon/4}/2}.
 \end{aligned}
\]
The first line uses the two-free estimate, and the second the one-free estimate; the finite union over interactions contributes \(C_u\). Both exponents tend to minus infinity faster than \(\log n\). On the remaining range the weight is bounded by \(e^{2^u n^{-.03}}\). The same statements hold for every subtuple of this fixed length.

\paragraph{The centered expansion.}
Choose \(L=L(P)\) sufficiently large and then a sufficiently large even \(u\). The moment \(\mathbb E(Z-1)^u\) is the integral of
\[
 \Phi_u(x_1,\ldots,x_u)
 =\sum_{I\subseteq[u]}(-1)^{u-|I|}
   \left(\int\prod_{j\in I}(1+a_{x_j})\,d\pi\right)^d
\]
against \(\sigma^{\otimes u}\). On \(w\le n^{-1-.03}\), expand each of the \(d\) factors into its nonempty interactions and the constant 1. A term selects interactions from distinct factors. If their union is \(U_0\), its coefficient in the alternating sum is
\(\sum_{I\supseteq U_0}(-1)^{u-|I|}\),
which vanishes unless \(U_0=[u]\). Thus every surviving list covers all tuple coordinates.

Lists with more than \(L\) interactions have total absolute contribution bounded by the geometric tail in \(2^u n^{-.03}\). Choose \(L\) so that this is smaller than the negative power needed for the final Markov bound. For a list of at most \(L\) interactions, some selected interaction contains at least \(u/L\) coordinates. Bound its integral by the mean-interaction estimate and the other factors by their small-set bounds. The total is at most
\[
 O_u(n^{L-.4u/L}).
\]
Choosing the even \(u\) with \(.4u/L>L+P+O(1)\) makes this small enough as well. On \(n^{-1-.03}<w\le n^{-\upsilon}\), bound each positive product in \(\Phi_u\) by \(e^{2^unw}\); the preceding exponential-weight estimate gives superpolynomial decay without cancellation.

\paragraph{Removing labels involved in large interactions.}
It remains to integrate tuples with \(w>n^{-\upsilon}\). From such a tuple choose a maximal retained set of coordinates \(R\) whose interactions all have magnitude at most \(n^{-\upsilon}\). Fix one possible \(R\) first, and later sum over the finitely many choices. Once the retained labels are fixed, maximality implies that every omitted label lies in a set \(B_R\) of extensions causing a large interaction with them. The count just proved gives
\[
 |B_R|\le2^u e^{n^{.03}}.
\]
This is a necessary condition on each omitted coordinate, not a claim of independence after conditioning on maximality. Under the original product law the separate necessary tests can be integrated independently, because \(B_R\) depends only on the retained coordinates.

For each positive product term in \(\Phi_u\), an omitted label involved in that term costs at most \(D_x^{-d}\) pointwise. We may use the same upper bound when it is not involved. Combining this factor with its necessary extension condition gives a charge at most
\[
 \begin{aligned}
 |B_R|\max_x\sigma(x)D_x^{-d}
 &\le\frac{2^n}{N}
       \exp\bigl(-.5gh+O(n^{.05})+n^{.03}+O_u(1)\bigr)\\
 &=n^{-\omega(1)}.
 \end{aligned}
\]
Here \(D_x^{-d}\le2^d e^{O(nb_*)}\), \(gh=n^{.09+o(1)}\), and \(2^n/N\le1\) eventually. The retained positive product has bounded integral by the exponential-weight estimate. At least one coordinate was omitted, so summing over retained sets and positive terms still gives a superpolynomially small contribution.

All parts of \(\mathbb E(Z-1)^u\) are therefore smaller than any prescribed negative power, after the stated choices of \(L,u\). Since \(Z<1/2\) implies \(|Z-1|^u>2^{-u}\), Markov proves \eqref{eq:source-8}.

\end{proof}

\subsection{Taking the assignment}
\label{sec:taking-the-assignment-jump}

\paragraph{Reference failure probability.}
Choose the slice tags \(i(s)\) independently from the compatible profile \(p\), draw the independent tuples at each tag's fixed \(y_0\), and use independent odd-row sampling conditional on those tuples. Denote the odd outputs by \(Y_b\). At an even \(v\), define the unnormalized row
\[
 M_v(x)=\sigma_v(x)\prod_{j\ {\rm outer}}
              \frac{1[x\sim_GY_{v^j}]}{D_x},
\]
where \(v^j\) is obtained from \(v\) by flipping outer bit \(j\); the row is zero off the indicated first support. Its mass is at least \(1/2\) except with arbitrarily strong polynomially small probability in this raw reference law. The internal failure probability for \(\sigma_v\) is already exponentially small. Given internal data with \(\sigma_v\) of mass 1, the outer neighbors lie in distinct slices. Their unconditional laws are independent copies of \(\pi\), so Lemma~\ref{lem:jump-outer-mass-tail} applies.

\paragraph{Avoidance at the tag stage.}
Fix a large constant \(P\), exceeding all needed polynomial dependency exponents and the requirements of Lemma~\ref{lem:clock-sampling} with \(B=4\). At each outer word \(s\), impose two conditions on the product tag law:
\begin{itemize}
\item the conditional raw probability that \(\|M_v\|_1<1/2\), for a given even role in this slice, is at most \(n^{-2P}\);
\item for every \(x\) on the first support of its tag, at least half the outer neighboring tags satisfy \(d_G(x;\pi_{i(s^j)})\le.8\).
\end{itemize}
The first conditional probability is identical for all even roles of the slice by inner translation. It reads the tags in the radius-one outer neighborhood. Markov applied to a stronger raw bound from the preceding paragraph makes its failure probability at most \(n^{-2P}\). For the second condition, compatibility gives exceptional tag probability at most \(\eta\) for every such \(x\). Independence of adjacent tags and a binomial union bound give total failure probability at most
\[
 N(4\eta)^{(n-h)/2}=e^{-\Omega(n)}.
\]
The tag events have dependency degree \(O(n^2)\). Lemma~\ref{lem:conditional-avoidance}, with charges of order \(n^{-2P}\), therefore supplies a law avoiding both kinds of failure.

We next choose a typical outcome of this tag-avoidance law. Write \(d=n-h\), and for a fixed host label define the normalized comparison means
\[
 A_s(y)=N\pi_{i(s)}(y),\qquad
 B_s(x)=N\alpha_{i(s)}(x)
       \prod_{j\le d}\frac{d_G(x;\pi_{i(s^j)})}{D_x}.
\]
The second expression is zero where its first-side row vanishes. These weights have cap at most \(2^ne^{-cn}\) on the tag gates for a fixed \(c>0\). For \(A_s\) use the width \(.02n\). For \(B_s\), use \(N\alpha_i\le e^{h\log2}\), \(D_x^{-d}\le2^de^{O(n^{.05})}\), and the second gate, which contributes at most \(.8^{d/2}\) in the numerator.

Apply scattered moments on the outer words, with words at bounded distance declared near. Their near fraction is \(n^{O(1)}2^{-d}\), and the product of this fraction, \(n\), and the cap tends to zero. At up to \(n\) separated neighborhoods remove the touching tag constraints; there are \(O(n^3)\) of them, with total charge \(o(1)\). In the restored product law,
\[
 \mathbb E A_s\le K,\qquad \mathbb E B_s\le K.
\]
For the second inequality, condition on \(i(s)\); each independent neighboring tag integrates its ratio to 1 on the compatible support, and then \eqref{eq:source-7} bounds \(\mathbb E_iN\alpha_i\). Thus, with high probability, both comparison means have bounded averages simultaneously for all labels. Fix tags having these bounds and all tag gates.

\paragraph{Avoidance at the tuple stage.}
At these fixed tags, the raw \(W\)-law is a product of the laws \(\rho_{y_0}^{\otimes k}\) at the even roles. Require at every even star that its conditional mass-failure probability under independent odd-row sampling be at most \(n^{-P}\). The tag gate and Markov make each tuple-stage bad event have probability at most \(n^{-P}\). It reads \(W\)-variables within full-cube distance 2, so the dependency degree is \(O(n^4)\). The conditional avoidance lemma applies again, now with charges of order \(n^{-P}\).

For a fixed second label \(y\), apply scattered moments to
\(Z_b(y)=Np_b^W(y)\)
under this tuple-avoidance law. The cap is \(e^{.02n}\), and a sufficiently large fixed full-cube neighborhood has fraction \(n^{O(1)}2^{-n}\). Their repeat cost is \(o(1)\). At separated rows remove the touching tuple constraints; at most \(O(n^5)\) constraints are removed for \(n\) queried neighborhoods, with charge sum \(o(1)\). Their raw means are \(A_{s(b)}(y)\), whose average is bounded by the selected tags. The resulting actual column sums are \(O(2^n/N)=o(1)\), simultaneously with high probability.

\paragraph{The odd injection.}
At a tuple history with these column bounds, apply the clock lemma, retaining this entering success indicator in subsequent comparisons. Its label count is \(N\), with \(\log N=O(n)\); column sums are eventually below \(\theta_0\); and row atoms are at most \(e^{.02n}/N<n^{-A_*}\). A mass-failure predicate reads the \(n\) odd neighbors of one even role, and each odd row occurs in \(n\) such predicates. Their product-law probabilities are at most \(n^{-P}\), as required by the chosen \(P\). The clock law therefore gives an odd injection avoiding every mass failure and joint upper comparison with the independent odd-row law.

\paragraph{Even loads and normalization.}
The normalized weights \(NM_v(x)\) have cap
\[
 N M_v(x)\le2^n\exp(-.5gh+O(n^{.05}))
               \le2^ne^{-.4gh}.
\]
A fixed-radius full-cube near relation again has fraction \(n^{O(1)}2^{-n}\), so its repeat cost is \(o(1)\), because \(gh=n^{.09+o(1)}\).

For up to \(n\) separated even stars, retain the successful tuple-history indicator and apply the clock comparison to their at most \(n^2\) distinct odd neighbors. This replaces the injection outputs by independent odd-row draws at fixed \(W\). We may then discard the nonlocal column-load indicator for an upper bound. Each remaining local row and weight reads only the radius-two \(W\)-scope of its star. Removing the avoidance constraints touching those scopes costs \(1+o(1)\) and restores independent raw tuple fields in the separated domains.

The resulting product means are exactly \(B_{s(v)}(x)\): the internal weight integrates to \(\alpha_{i(s(v))}\), and each outer label contributes its degree into the corresponding \(\pi_{i(s^j)}\). Their bounded averages give, by scattered moments, simultaneous column sums
\[
 \sum_vM_v(x)=O(2^n/N)=o(1),
\]
hence at most \(1/2\) with high probability on the entering successes. This uses only \(N/2^n\to\infty\), with no rate assumption. Clock avoidance has ensured that every row has mass at least \(1/2\). Normalizing each row thus increases column sums by at most 2. The resulting probability rows lie on common neighbors of the odd injection and have column sums at most 1, so Hall's theorem completes the cube. This excludes the linear-budget jump.

\end{proof}

\begin{corollary}[The remaining deep regime]\label{cor:remaining-deep-regime}
The possibility remaining in this analysis has \textbf{both} \(H_L^\dagger\ge1\), i.e., for each fixed error exponent slack \(\epsilon>0\) one has deep discrepancy \(\le n^{-1+\epsilon}\) on retained labels at some budgets \(n^x,\alpha n\) with constants \(x,\alpha>0\), in both orientations.
\end{corollary}

\begin{proof}
Proposition~\ref{prop:intermediate-bias-exclusion} leaves both \(H^\dagger\ge1\) and excludes \(0<H_L^\dagger<1\). Proposition~\ref{prop:linear-budget-jump-exclusion} also excludes \(H_L^\dagger=0\); since these exponents are nonnegative, both linear limits are at least 1. Given \(0<\epsilon<1\), choose a rational \(h_0\) with \(1-\epsilon<h_0<1\), and then positive rational budgets with \(H_L(x,\alpha)>h_0\) in each orientation. Stabilized absence at \(h_0\) gives discrepancy less than \(n^{-h_0}\le n^{-1+\epsilon}\). Larger error slacks follow by weakening such a bound.
\end{proof}
